\chapter{Klasyfikacja topologicznych rozmaitości czterowymiarowych}
\label{ch:freedman-klasyfikacja}
\index{Freedman!klasyfikacja wymiaru cztery}

W Rozdziale~\ref{ch:freedman-formy} obliczaliśmy formy
przecięcia i oglądaliśmy przeszkodę dla gładkiego
triku Whitneya. Uchwyty Cassona, opisane w
Rozdziale~\ref{ch:uchwyty-cassona}, stanowią drogę
do topologicznego twierdzenia o zanurzaniu dysku.
Poniżej podamy precyzyjny wynik klasyfikacyjny
Freedmana jako wejście zewnętrzne i wyprowadzimy
z niego wnioski. Każde stwierdzenie o homeomorfizmie
dotyczy kategorii TOP; typy gładkie wymagają
osobnych argumentów.

\section{Niezmiennik Kirby'ego--Siebenmanna}
\label{sec:freedman-ks}
\index{niezmiennik Kirby'ego--Siebenmanna}

Niech $X$ będzie zamkniętą, spójną, topologiczną
rozmaitością $4$-wymiarową. Jej stabilna wiązka
styczna jest w kategorii TOP mikrowiązką. Teoria
Kirby'ego--Siebenmanna przypisuje jej klasę
\begin{equation}\label{eq:freedman-ks}
 \operatorname{ks}(X)\in H^4(X;\Z/2).
\end{equation}
Po ewaluacji na mod $2$ klasie fundamentalnej
otrzymujemy liczbę w $\Z/2$, również oznaczaną
$\operatorname{ks}(X)$. Intuicyjnie mierzy ona
przeszkodę do \emph{stabilnego} wygładzenia: dla
zamkniętej rozmaitości $4$-wymiarowej
$\operatorname{ks}(X)=0$ wtedy i tylko wtedy,
gdy $X\times\R$ dopuszcza strukturę gładką.
Jest to wynik zewnętrzny teorii wygładzania,
\href{https://archive.mpim-bonn.mpg.de/4789/2/FreedmanQuinn-TopologyOf4Manifolds-Reformatted2013.pdf}
{Freedman--Quinn, \S\,8.7--8.8 i \S\,10.2.2},
oparty na
\href{https://www.maths.ed.ac.uk/~v1ranick/papers/ks.pdf}
{Kirby--Siebenmann, \emph{Foundational Essays on
Topological Manifolds, Smoothings, and Triangulations}
(1977)}.

W szczególności struktura gładka lub PL na $X$
wymusza $\operatorname{ks}(X)=0$. W wymiarze cztery
odwrotnej implikacji \emph{nie} wolno wyciągać z
tego twierdzenia: zerowanie przeszkody stabilnej
nie buduje automatycznie wygładzenia samego $X$.
Niezmiennik jest zachowywany przez homeomorfizmy,
a dla sum spójnych zachodzi
\[
 \operatorname{ks}(X\#Y)
   =\operatorname{ks}(X)+\operatorname{ks}(Y)
       \quad\text{w }\Z/2.
\]
Te własności również przyjmujemy z powyższego źródła.

Jeżeli $X$ jest dodatkowo jednospójna i zorientowana,
parzystość $Q_X$ odpowiada warunkowi
$w_2(X)=0$: jest to wzór Wu
$Q_X(x,x)\equiv\langle w_2(X),x\rangle\pmod2$
wraz z brakiem $2$-torsji w $H_1(X)$.
Jest to topologiczna wersja warunku spinowego;
również ją przyjmujemy jako wynik zewnętrzny,
\href{https://archive.mpim-bonn.mpg.de/4789/2/FreedmanQuinn-TopologyOf4Manifolds-Reformatted2013.pdf}
{Freedman--Quinn, \S\,10.2.1--10.2.2}.
Bez hipotezy o $H_1$ sama parzystość formy na
całkowitej homologii nie musi wykrywać całego $w_2$.

\section{Twierdzenie Freedmana jako wynik zewnętrzny}
\label{sec:freedman-classification}

\begin{twierdzenie}[Freedman: klasyfikacja jednospójna;
wynik zewnętrzny]
\label{thm:freedman-classification}
Niech $Q$ będzie symetryczną unimodularną formą
całkowitą na wolnej grupie abelowej skończonej rangi,
a $k\in\Z/2$. Jeśli $Q$ jest parzysta, załóżmy
dodatkowo
\begin{equation}\label{eq:freedman-even-ks}
 k\equiv\frac{\sigma(Q)}8\pmod2.
\end{equation}
Wtedy istnieje zamknięta, spójna, zorientowana,
jednospójna \emph{topologiczna} rozmaitość
$X^4$ z $Q_X\cong Q$ i $\operatorname{ks}(X)=k$.
Jeśli $Q$ jest nieparzysta, obie wartości $k=0,1$
są dopuszczalne.

Ponadto dwie takie rozmaitości są homeomorficzne
z zachowaniem orientacji wtedy i tylko wtedy,
gdy ich formy przecięcia są izometryczne i mają
ten sam niezmiennik Kirby'ego--Siebenmanna.
W wersji z ustaloną izometrią form homeomorfizm
można wybrać tak, by indukował właśnie tę izometrię.
\end{twierdzenie}

Wzór~\eqref{eq:freedman-even-ks} ma sens, ponieważ
sygnatura parzystej formy unimodularnej jest podzielna
przez $8$; jak zaznaczono po
Stwierdzeniu~\ref{prop:freedman-E8-properties},
ogólny fakt algebraiczny przyjmujemy z literatury.
Dla jednospójnej $X$ parzystość $Q_X$ oznacza także
$w_2(X)=0$, więc $X$ jest spinowa. W tym przypadku
zewnętrzna formuła dla zamkniętych topologicznych
rozmaitości spinowych brzmi
$\operatorname{ks}(X)\equiv\sigma(X)/8\pmod2$;
zob. Freedman--Quinn, \S\,10.2.2. Wyjaśnia ona
konieczność warunku~\eqref{eq:freedman-even-ks},
natomiast część istnienia Twierdzenia~\ref{thm:freedman-classification}
zapewnia jego wystarczalność.
Oryginalne Twierdzenie 1.5 w pracy
\href{https://www.maths.gla.ac.uk/~mpowell/1982_The%20topology%20of%20four-dimensional%20manifolds.pdf}
{M.\ H.\ Freedman, \emph{The topology of four-dimensional
manifolds}, J. Differential Geom. 17 (1982), 357--453}
jest sformułowane dla rozmaitości gładkich poza jednym
punktem. Podaną wyżej pełną postać w kategorii TOP,
razem z warunkiem na $\operatorname{ks}$, przyjmujemy z
\href{https://archive.mpim-bonn.mpg.de/4789/2/FreedmanQuinn-TopologyOf4Manifolds-Reformatted2013.pdf}
{M.\ H.\ Freedman i F.\ Quinn, \emph{Topology of
4-Manifolds}, \S\,10.1}.
W źródłach dowód korzysta z topologicznego zanurzania
dysków i uchwytów Cassona. Poprzednie sekcje tego
rozdziału nie dowodzą istnienia ani jednoznaczności
homeomorfizmu z twierdzenia.

\section{Dwie konstrukcje w dowodzie Freedmana}
\label{sec:freedman-constructive-examples}
\index{Freedman!konstrukcja rozmaitości o formie $E_8$}

Twierdzenie~\ref{thm:freedman-classification} jest
wynikiem zewnętrznym, ale jego pierwszą konstrukcję
możemy prześledzić na obliczonym już plumbowaniu.
Oddzielimy rachunki homologii i formy od istnienia
topologicznego wypełnienia brzegu. Drugą konstrukcją
będzie rozmaitość o formie $[1]$, której nie da się
odróżnić od $\C P^2$ samą formą przecięcia.

\begin{przyklad}[Od plumbowania do $X_{E_8}$]
\label{prz:freedman-E8-gluing}
Niech $P$ będzie plumbowaniem ośmiu zorientowanych
wiązek $D^2\to P_i\to S^2$ o liczbie Eulera $+2$.
Ich rdzenie są sferami o samoprzecięciu $+2$;
plumbowania wykonujemy według grafu z
Rysunku~\ref{fig:freedman-E8-graph}, wybierając
znaki przecięć sąsiednich rdzeni równe $-1$.
Równoważnie $P$ otrzymujemy z $D^4$ przez osiem
uchwytów $2$ o macierzy sprzężeń $E_8$.
Nie ma uchwytów $1$, zatem $\pi_1(P)=1$, a
Stwierdzenie~\ref{prop:freedman-link-matrix} daje
\[
 H_2(P;\Z)\cong\Z^8,\qquad Q_P=E_8,
 \qquad H_1(\partial P;\Z)=H_2(\partial P;\Z)=0.
\]
Ostatnia równość korzysta z
$\det E_8=1$ ze Stwierdzenia~\ref{prop:freedman-E8-properties}.
Brzeg $\Sigma=\partial P$ jest więc całkowitą
sferą homologiczną. Według Freedmana--Quinna jest
to sfera homologiczna Poincarégo z odpowiednią
orientacją; sam rachunek homologii nie identyfikuje
jej jako $S^3$.

Głęboki krok pochodzi z twierdzenia Freedmana:
\emph{każda} całkowita sfera homologiczna wymiaru
trzy ogranicza kontraktowalną topologiczną
rozmaitość $4$-wymiarową. Wybierzmy takie $C$
z $\partial C\cong-\Sigma$ i sklejmy
\[
 X_{E_8}=P\cup_{\Sigma}C.
\]
Istnienie $C$ jest wynikiem zewnętrznym,
\href{https://www.maths.gla.ac.uk/~mpowell/1982_The%20topology%20of%20four-dimensional%20manifolds.pdf}
{Freedman, \emph{The topology of four-dimensional
manifolds}, twierdzenie $1.4'$}; zob. też
\href{https://archive.mpim-bonn.mpg.de/4789/2/FreedmanQuinn-TopologyOf4Manifolds-Reformatted2013.pdf}
{Freedman--Quinn, wniosek 9.7.2 i \S\,10.5.1}.

Pozostałe kroki są już rachunkami. Twierdzenie
van Kampena opisuje $\pi_1(X_{E_8})$ jako iloraz
$\pi_1(P)*\pi_1(C)$ przez relacje pochodzące od
$\pi_1(\Sigma)$. Oba czynniki są trywialne, bo
$P$ jest jednospójne, a $C$ kontraktowalne;
stąd $\pi_1(X_{E_8})=1$.
Fragment ciągu Mayera--Vietorisa ma postać
\[
 0=H_2(\Sigma;\Z)\longrightarrow
 H_2(P;\Z)\oplus H_2(C;\Z)
 \longrightarrow H_2(X_{E_8};\Z)
 \longrightarrow H_1(\Sigma;\Z)=0.
\]
Ponieważ $H_2(C)=0$, inkluzja $P\hookrightarrow X_{E_8}$
indukuje izomorfizm na $H_2$. Sfery zerowe
plumbowanych wiązek reprezentują bazę $H_2(P)$
i przecinają się we wnętrzu $P$, więc po sklejeniu
zachowują macierz $E_8$. Otrzymaliśmy zatem
zamkniętą, jednospójną rozmaitość topologiczną
o $Q_{X_{E_8}}=E_8$ i $\sigma(X_{E_8})=8$.
Forma jest parzysta, więc $X_{E_8}$ jest spinowa;
zewnętrzny wzór z Sekcji~\ref{sec:freedman-ks}
daje $\operatorname{ks}(X_{E_8})=8/8=1\pmod2$.
To obliczenie wyklucza wygładzenie $X_{E_8}$.
Nie dostarcza ono samodzielnej konstrukcji $C$.
\end{przyklad}

\begin{przyklad}[Drugi typ o formie nieparzystej rzędu jeden]
\label{prz:freedman-odd-twin}
\index{Freedman!drugi typ o formie $[1]$}
Rozważmy $Z=X_{E_8}\#\overline{\C P^2}$. Addytywność
formy i niezmiennika Kirby'ego--Siebenmanna daje
\[
 Q_Z=E_8\oplus[-1],\qquad
 \operatorname{ks}(Z)=1.
\]
Forma $Q_Z$ ma rząd $9$, sygnaturę $7$, jest
nieparzysta i nieokreślona. Klasyfikacja całkowitych
nieparzystych form nieokreślonych daje izometrię
\[
 E_8\oplus[-1]\cong8[1]\oplus[-1].
\]
Jest to twierdzenie algebraiczne cytowane przez
Freedmana--Quinna w \S\,10.2.1; jawny rachunek
odpowiedniej bazy, po odwróceniu znaków, znajduje się
również w Rozdziale~\ref{ch:egzotyczne-r4}.

Wybierzmy w tej izometrii podformę
$7[1]\oplus[-1]$, realizowaną przez gładką
rozmaitość $M=7\C P^2\#\overline{\C P^2}$.
Jej dopełnienie ortogonalne ma formę $[1]$,
na której $w_2$ jest niezerowa, bo kwadrat
generatora jest nieparzysty. Freedman--Quinn
sprawdzają w tym jednospójnym przypadku
pozostałe warunki dotyczące przecięć
i samoprzecięć. Ich zewnętrzne
\emph{twierdzenie o rozkładzie sumy spójnej}
pozwala zrealizować tę podformę jako rozkład
\[
 Z\cong M\#N,\qquad Q_N=[1].
\]
Właśnie ten krok, a nie sama izometria kratek,
korzysta z głębokiej geometrii wymiaru cztery;
zob. \href{https://archive.mpim-bonn.mpg.de/4789/2/FreedmanQuinn-TopologyOf4Manifolds-Reformatted2013.pdf}
{Freedman--Quinn, twierdzenie 10.4 i \S\,10.5.1
w numeracji wersji z 2013 r.}; w wydaniu z 1990 r.
są to twierdzenie 10.3 i dowód twierdzenia 10.1
w \S\,10.4.
Z addytywności otrzymujemy
$\operatorname{ks}(N)=\operatorname{ks}(Z)-
\operatorname{ks}(M)=1$, gdyż $M$ jest gładka.
Ponadto $N$ jest jednospójna na mocy twierdzenia
van Kampena dla sumy spójnej. Ma tę samą formę
$[1]$ co $\C P^2$, lecz inne $\operatorname{ks}$;
stąd nie jest z nią homeomorficzna i nie może
być gładka. Istnienie obu typów w nieparzystym
przypadku Twierdzenia~\ref{thm:freedman-classification}
ma więc konkretny punkt wyjścia w $X_{E_8}$.
\end{przyklad}

\section{Jak zanurzanie dysków wchodzi do klasyfikacji}
\label{sec:freedman-route}

W przypadku jednospójnym grupa podstawowa jest
trywialna, więc spełnia warunek
Twierdzenia~\ref{thm:casson-disk-embedding}.
Twierdzenie to zastępuje odpowiednio przygotowane
immersyjne dyski lokalnie płaskimi zanurzeniami, lecz
\emph{nie} klasyfikuje jeszcze rozmaitości: nie
konstruuje samo formy przecięcia, nie wyznacza
niezmiennika Kirby'ego--Siebenmanna i nie daje
homeomorfizmu między dwiema zadanymi rozmaitościami.

Pierwszym dalszym krokiem jest przejście od dysków
do sfer. W używanej tutaj szczególnej postaci
zewnętrzne \emph{twierdzenie o zanurzaniu sfer} mówi:
jeżeli w spójnej, zorientowanej topologicznej
rozmaitości $M^4$ z $\pi_1(M)=1$ mamy rodzinę obramowanych,
immersyjnych sfer $A_i$ z
$\lambda(A_i,A_j)=0$ dla $i\ne j$ i
$\mu(A_i)=0$, a także obramowane immersyjne sfery
algebraicznie dualne do $A_i$, to rodzinę $A_i$ można
przez homotopię regularną doprowadzić do parami
rozłącznych, lokalnie płasko zanurzonych obramowanych
sfer mających geometrycznie dualne immersyjne sfery.
W tej obramowanej postaci normalna liczba Eulera
$e(A_i)$ wynosi zero. Ponieważ $\pi_1(M)=1$,
to $\lambda(A_i,A_i)=2\mu(A_i)+e(A_i)=0$;
spełniony jest zatem warunek homotopii regularnej,
a zanurzone sfery mają trywialne wiązki normalne.
Jest to szczególny
przypadek \href{https://arxiv.org/pdf/2006.05209}
{Powell--Ray--Teichner, \emph{The 4-dimensional disc
embedding theorem and dual spheres} (2025),
Twierdzenie B i Uwaga 2.2}. W ogólnej postaci twierdzenie
rozróżnia zredukowaną liczbę samoprzecięć od
liczby Eulera normalnego i wymaga $w_2=0$ na
klasach sfer dualnych. Podane tutaj obramowanie
spełnia ten ostatni warunek.

Geometryczna dualność pozwala kontrolować
chirurgię na zanurzonych sferach. Dla jednej
obramowanej sfery $S$ usuwamy
$S^2\times\operatorname{int}D^2$ i wklejamy
$D^3\times S^1$ wzdłuż wspólnego
$S^2\times S^1$. Oznaczmy przez $C$ domknięcie
dopełnienia usuniętego otoczenia, a przez $m$
meridian $S$ w $\partial C$. Twierdzenie van
Kampena daje
\[
 \pi_1(M)\cong\pi_1(C)/\langle\!\langle m\rangle\!\rangle,
 \qquad
 \pi_1(M_{\mathrm{chir}})\cong\pi_1(C),
\]
ponieważ inkluzja $S^2\times S^1\hookrightarrow
D^3\times S^1$ indukuje izomorfizm grup
podstawowych. Gdy $S$ ma
geometrycznie dualną sferę, meridian ogranicza
w $C$ immersyjny dysk, dokładnie jak
w dowodzie Stwierdzenia~\ref{prop:casson-complement-pi1}.
Zatem $m=1$ w $\pi_1(C)$ i chirurgia zachowuje
grupę podstawową.
Jest to jeden z konkretnych powodów, dla których
wniosek o sferach dualnych w twierdzeniu o
zanurzaniu nie jest tylko dodatkiem do jego tezy.

Następne kroki są \emph{zewnętrzne}: algebraiczna
teoria chirurgii organizuje jądra przekształceń
normalnych w pary sfer do zanurzenia; topologiczna
chirurgia w wymiarze cztery używa tych zanurzeń
do usunięcia jądra; twierdzenie o $s$-kobordyzmie
w wymiarze pięć dla dobrej grupy podstawowej
(tu trywialnej) identyfikuje odpowiedni topologiczny
$s$-kobordyzm, którego inkluzje końców są prostymi
równoważnościami homotopijnymi, z produktem.
Dodatkowe argumenty istnienia rozmaitości
o zadanej formie oraz rachunek niezmiennika
Kirby'ego--Siebenmanna domykają klasyfikację.
Dokładne twierdzenia i ich hipotezy znajdują się
u \href{https://archive.mpim-bonn.mpg.de/4789/2/FreedmanQuinn-TopologyOf4Manifolds-Reformatted2013.pdf}
{Freedmana--Quinna, rozdziały 7, 10 i 11},
zwłaszcza Twierdzenie 7.2 oraz \S\,10.1 i \S\,11.6.
Powiązanie tych narzędzi wyjaśnia rolę
Twierdzenia~\ref{thm:casson-disk-embedding},
ale nie jest dowodem
Twierdzenia~\ref{thm:freedman-classification}.

Rysunek~\ref{fig:freedman-classification-route} zbiera
kolejność zależności dla przypadku $\pi_1=1$.
Sfery algebraicznie dualne są częścią danych wejściowych
twierdzenia o zanurzaniu; geometrycznie dualne sfery
są jego wnioskiem. Strzałki oznaczają użycie wyników
zewnętrznych z hipotezami podanymi wyżej, a nie
samodzielną konstrukcję klasyfikacji.

\begin{figure}[htbp]
\centering
\begin{tikzpicture}[
  font=\small,>=Latex,
  routebox/.style={draw=manifoldblue!75!black,
    fill=manifoldblue!7,rounded corners=2pt,
    align=center,text width=3.65cm,minimum height=1.25cm,
    inner sep=4pt},
  routearrow/.style={->,thick,manifoldblue!80!black}]
 \node[routebox] (immersed) at (-5.05,0)
   {immersyjne dyski i sfery; dualność algebraiczna};
 \node[routebox] (towers) at (0,0)
   {wieże Cassona i uogólnienia Freedmana--Quinna};
 \node[routebox] (embedded) at (5.05,0)
   {lokalnie płaskie dyski i sfery; dualność geometryczna};
 \node[routebox] (surgery) at (5.05,-2.1)
   {chirurgia TOP i kontrola $\pi_1$};
 \node[routebox] (cobordism) at (0,-2.1)
   {topologiczny $s$-kobordyzm w wymiarze pięć};
 \node[routebox] (classification) at (-5.05,-2.1)
   {klasyfikacja przez $(Q,\operatorname{ks})$};
 \draw[routearrow] (immersed)--(towers);
 \draw[routearrow] (towers)--(embedded);
 \draw[routearrow] (embedded)--(surgery);
 \draw[routearrow] (surgery)--(cobordism);
 \draw[routearrow] (cobordism)--(classification);
\end{tikzpicture}
\caption{Schemat zależności w przypadku jednospójnym.
Strzałki wymagają zaniku przeszkód algebraicznych i pozostałych
hipotez odpowiednich twierdzeń. Twierdzenia o zanurzaniu
dysków i sfer oraz o topologicznej
chirurgii i $s$-kobordyzmie są wynikami zewnętrznymi.
Realizacja form i rachunek $\operatorname{ks}$ są dodatkowymi
wejściami do ostatniego kroku.}
\label{fig:freedman-classification-route}
\end{figure}

\noindent\textbf{Mechanizm jednoznaczności.}
Niech $M$ i $N$ będą zamkniętymi, zorientowanymi,
jednospójnymi rozmaitościami o izometrycznych
formach i równych niezmiennikach $\operatorname{ks}$.
Twierdzenie Hurewicza utożsamia tu $\pi_2(M)$
z $H_2(M)$, a $\pi_2(N)$ z $H_2(N)$.
W dowodzie Freedmana--Quinna zewnętrzne
twierdzenie o rozkładzie sumy spójnej realizuje
te dane geometrycznie: prowadzi do rozkładu
$N\cong M\#P$ albo do analogicznego rozkładu
wariantu $N$ o przeciwnym $\operatorname{ks}$
w przypadku formy nieparzystej. Dopełnienie $P$
jest jednospójne i ma zerowe $H_2$; dualność
Poincarégo czyni je sferą homologiczną wymiaru
cztery. Topologiczna hipoteza Poincarégo,
również wynik zewnętrzny Freedmana, daje
$P\cong S^4$. Równość $\operatorname{ks}(M)
=\operatorname{ks}(N)$ wybiera pierwszą możliwość;
dla formy parzystej wariant z przeciwnym
$\operatorname{ks}$ nie występuje. To wyjaśnia
rolę obu niezmienników, lecz nie zastępuje dowodu
twierdzenia o rozkładzie ani hipotezy Poincarégo;
zob. \href{https://archive.mpim-bonn.mpg.de/4789/2/FreedmanQuinn-TopologyOf4Manifolds-Reformatted2013.pdf}
{Freedman--Quinn, \S\,10.5.1 i wniosek 7.2.1
w numeracji wersji z 2013 r.}.

\begin{wniosek}[Topologiczna hipoteza Poincarégo w wymiarze cztery]
\label{cor:freedman-poincare-4}
Jeżeli zamknięta, spójna, jednospójna topologiczna
rozmaitość $X^4$ ma homologie całkowite sfery $S^4$,
to $X$ jest homeomorficzna z $S^4$.
W szczególności każda topologiczna sfera homotopijna
wymiaru cztery jest homeomorficzna z $S^4$.
\end{wniosek}
\begin{proof}
Ponieważ $H_4(X;\Z)\cong\Z$, wybieramy orientację.
Warunek homologiczny daje $H_2(X;\Z)=0$, więc
$Q_X$ jest formą zerową na grupie rangi zero.
Jest parzysta, unimodularna w sensie pustej macierzy
o wyznaczniku $1$, i ma sygnaturę $0$.
Z~\eqref{eq:freedman-even-ks} wynika
$\operatorname{ks}(X)=0$. Standardowa $S^4$
ma te same dane, więc część jednoznaczności
Twierdzenia~\ref{thm:freedman-classification}
daje homeomorfizm. Wniosek nie stwierdza
istnienia dyfeomorfizmu.
\end{proof}

\begin{przyklad}[Cztery dane topologiczne]
\label{prz:freedman-examples}
Forma $H$ jest parzysta, ma sygnaturę $0$ i
$\operatorname{ks}=0$, więc jej jedynym zorientowanym
typem topologicznym jest $S^2\times S^2$.
Forma $[1]$ jest nieparzysta: istnieją dokładnie
dwa typy. Jeden to $\C P^2$ z
$\operatorname{ks}=0$; drugi, z
Przykładu~\ref{prz:freedman-odd-twin}, ma
$\operatorname{ks}=1$ i nie może mieć struktury
gładkiej ani PL.

Forma $E_8$ jest parzysta i ma sygnaturę $8$.
Konstrukcja z Przykładu~\ref{prz:freedman-E8-gluing}
daje $X_{E_8}$, a Twierdzenie Freedmana jej
jedyny zorientowany typ topologiczny o tej formie;
a~\eqref{eq:freedman-even-ks} daje
$\operatorname{ks}(X_{E_8})=1$.
Wobec własności z Sekcji~\ref{sec:freedman-ks}
$X_{E_8}$ nie ma struktury gładkiej ani PL.
Natomiast forma $E_8\oplus E_8$ ma sygnaturę $16$,
więc odpowiadający jej typ topologiczny ma
$\operatorname{ks}=0$. Tego ostatniego rachunku
nie interpretujemy jako dowodu gładkości:
Wniosek~\ref{cor:donaldson-2e8} wykluczy
wygładzenie tej rozmaitości za pomocą
twierdzenia Donaldsona.
\end{przyklad}

\begin{table}[htbp]
\centering
\renewcommand{\arraystretch}{1.12}
\setlength{\tabcolsep}{9pt}
\begin{tabular}{@{}lccclc@{}}
\hline
Forma & Parzystość & $\sigma$ & $\operatorname{ks}$
      & Typ TOP & DIFF \\
\hline
$H$ & parzysta & $0$ & $0$ & $S^2\times S^2$ & tak \\
$[1]$ & nieparzysta & $1$ & $0$ & $\C P^2$ & tak \\
$[1]$ & nieparzysta & $1$ & $1$ & $X_{[1],1}$ & nie \\
$E_8$ & parzysta & $8$ & $1$ & $X_{E_8}$ & nie \\
$2E_8$ & parzysta & $16$ & $0$ & $X_{2E_8}$ & nie \\
\hline
\end{tabular}
\caption{Przykłady zamkniętych, spójnych, zorientowanych
i jednospójnych rozmaitości. Każdy wiersz wyznacza dokładnie
jeden typ TOP; $X_{[1],1}$ oznacza drugi typ o formie $[1]$.
W kolumnie DIFF „tak” oznacza istnienie co najmniej jednej
struktury gładkiej, nie jej jednoznaczność.
Brak wygładzenia dla $2E_8$ wynika z
Wniosku~\ref{cor:donaldson-2e8}, mimo że
$\operatorname{ks}=0$.}
\label{tab:freedman-examples}
\end{table}

Twierdzenie~\ref{thm:freedman-classification}
dotyczy tylko zamkniętych, zorientowanych,
jednospójnych rozmaitości w kategorii TOP.
Przypadek ogólnej grupy podstawowej wymaga
dodatkowych danych i hipotez. Własności gładkich
struktur na tej samej topologicznej rozmaitości
wymagają dalszych narzędzi, w szczególności
teorii Donaldsona; żaden rachunek samego
wyznacznika, parzystości i sygnatury ich nie rozstrzyga.

\par\smallskip
{\footnotesize\noindent\emph{Dalsza lektura:}
\href{https://www.maths.gla.ac.uk/~mpowell/1982_The%20topology%20of%20four-dimensional%20manifolds.pdf}
{M.\ H.\ Freedman, \emph{The topology of four-dimensional
manifolds}, J. Differential Geom. 17 (1982), 357--453},
oryginalna konstrukcja;
\href{https://archive.mpim-bonn.mpg.de/4789/2/FreedmanQuinn-TopologyOf4Manifolds-Reformatted2013.pdf}
{M.\ H.\ Freedman i F.\ Quinn, \emph{Topology of
4-Manifolds} (1990)}, zwłaszcza \S\,10.1--10.2:
pełna postać klasyfikacji i niezmiennik
Kirby'ego--Siebenmanna.\par}
